\documentclass[a4paper,11pt]{article}

% --- Dane na stronę tytułową ---------------------------------------------
\newcommand{\autor}{Eryk Hadała}
\newcommand{\grupa}{Grupa laboratoryjna 1}
\newcommand{\prowadzacy}{dr inż. Grzegorz Rogus}
\newcommand{\datadokumentu}{26 października 2026}

% --- Język i czcionki -----------------------------------------------------
\usepackage{fontspec}
\usepackage[polish]{babel}
\setmainfont{texgyrepagella}[
  Extension=.otf,
  UprightFont=*-regular,
  BoldFont=*-bold,
  ItalicFont=*-italic,
  BoldItalicFont=*-bolditalic]
\setsansfont{texgyreheros}[
  Extension=.otf,
  UprightFont=*-regular,
  BoldFont=*-bold,
  ItalicFont=*-italic,
  BoldItalicFont=*-bolditalic]
\setmonofont{lmmono10-regular.otf}[Scale=0.95]
\usepackage{csquotes}
\usepackage{microtype}

% --- Układ strony ---------------------------------------------------------
\usepackage[a4paper,top=2.5cm,bottom=2.5cm,left=2.3cm,right=2.3cm,headheight=22pt]{geometry}
\usepackage{xcolor}
\definecolor{accent}{HTML}{2F5D50}
\definecolor{accentlight}{HTML}{E8F0EC}
\definecolor{muted}{HTML}{5F6B66}
\definecolor{ok}{HTML}{2E7D32}
\definecolor{warn}{HTML}{B26A00}
\definecolor{bad}{HTML}{B3261E}

\usepackage{titlesec}
\titleformat{\part}[display]
  {\sffamily\bfseries\color{accent}\centering}
  {\Large\MakeUppercase{\partname}~\thepart}{0.6em}{\Huge}
\titleformat{\section}{\sffamily\Large\bfseries\color{accent}}{\thesection}{0.8em}{}
\titleformat{\subsection}{\sffamily\large\bfseries}{\thesubsection}{0.8em}{}
\titleformat{\subsubsection}{\sffamily\normalsize\bfseries\color{muted}}{\thesubsubsection}{0.8em}{}
\titlespacing*{\section}{0pt}{1.6em}{0.7em}

\usepackage{fancyhdr}
\pagestyle{fancy}
\fancyhf{}
\fancyhead[L]{\sffamily\small\color{muted}Listy zakupów — plan i raport z testów}
\fancyhead[R]{\sffamily\small\color{muted}\thepage}
\renewcommand{\headrulewidth}{0.4pt}
\renewcommand{\headrule}{\hbox to\headwidth{\color{accent!40}\leaders\hrule height \headrulewidth\hfill}}

\usepackage{enumitem}
\setlist{itemsep=0.25em,topsep=0.4em}
\setlength{\parindent}{0pt}
\setlength{\parskip}{0.55em}

% --- Tabele, wykresy, ramki -----------------------------------------------
\usepackage{array}
\usepackage{booktabs}
\usepackage{tabularx}
\usepackage{xltabular}
\usepackage{multirow}
\usepackage{colortbl}
\usepackage{float}
\usepackage{caption}
\captionsetup{font=small,labelfont={sf,bf,color=accent},skip=6pt}
\renewcommand{\arraystretch}{1.25}
\newcolumntype{L}{>{\raggedright\arraybackslash}X}
\newcolumntype{C}{>{\centering\arraybackslash}X}

\usepackage{pgfplots}
\pgfplotsset{compat=1.18,
  /pgf/number format/use comma,
  /pgf/number format/1000 sep={\,},
  legend style={/tikz/every even column/.append style={column sep=10pt}},
  every axis legend/.append style={font=\small}}

\usepackage[most]{tcolorbox}
\newtcolorbox{werdykt}[1][]{colback=accentlight,colframe=accent,boxrule=0.6pt,
  arc=2pt,left=8pt,right=8pt,top=6pt,bottom=6pt,fonttitle=\sffamily\bfseries,title={#1}}
\newtcolorbox{uwaga}{colback=white,colframe=warn,boxrule=0.6pt,arc=2pt,
  left=8pt,right=8pt,top=5pt,bottom=5pt}

\usepackage[hidelinks]{hyperref}
\hypersetup{colorlinks=true,linkcolor=accent,urlcolor=accent,
  pdftitle={Listy zakupów — plan testów i raport z testów},pdfauthor={\autor}}

% --- Skróty ---------------------------------------------------------------
\newcommand{\code}[1]{\texttt{\small #1}}
\newcommand{\zal}{\textcolor{ok}{\textbf{zaliczony}}}
\newcommand{\pass}{\textcolor{ok}{\textbf{OK}}}
\newcommand{\sev}[1]{\textsf{\small\bfseries #1}}
\newcommand{\wysoki}{\textcolor{bad}{\sev{wysoki}}}
\newcommand{\sredni}{\textcolor{warn}{\sev{średni}}}
\newcommand{\niski}{\textcolor{accent}{\sev{niski}}}
\newcommand{\info}{\textcolor{muted}{\sev{informacyjny}}}
\newcommand{\plik}[1]{\multicolumn{4}{@{}l}{\cellcolor{accentlight}\textsf{\footnotesize\color{accent} #1}}\\}

% Tabela przypadków testowych: ID | przypadek | dane | oczekiwany wynik
\newenvironment{przypadki}
  {\footnotesize
   \xltabular{\textwidth}{@{}>{\sffamily\bfseries\color{accent}}p{1.15cm} L L L@{}}
   \toprule
   \textcolor{black}{ID} & \textsf{\textbf{Przypadek}} & \textsf{\textbf{Dane wejściowe}} & \textsf{\textbf{Oczekiwany wynik}}\\
   \midrule
   \endfirsthead
   \toprule
   \textcolor{black}{ID} & \textsf{\textbf{Przypadek}} & \textsf{\textbf{Dane wejściowe}} & \textsf{\textbf{Oczekiwany wynik}}\\
   \midrule
   \endhead
   \bottomrule
   \endfoot}
  {\endxltabular}

\begin{document}

% =========================================================================
\hypersetup{pageanchor=false}
\begin{titlepage}
  \sffamily
  \vspace*{2.5cm}
  {\color{accent}\rule{\linewidth}{1.2pt}}\par\vspace{0.8cm}
  {\large\color{muted} Laboratorium Testowania Oprogramowania}\par\vspace{0.5cm}
  {\fontsize{30}{36}\selectfont\bfseries\color{accent} Listy zakupów}\par\vspace{0.4cm}
  {\LARGE Plan testów i raport z testów}\par\vspace{0.8cm}
  {\color{accent}\rule{\linewidth}{1.2pt}}\par
  \vspace{1.2cm}
  \begin{tabular}{@{}ll@{}}
    \textcolor{muted}{Autor:} & \autor \\
    \textcolor{muted}{Grupa:} & \grupa \\
    \textcolor{muted}{Prowadzący:} & \prowadzacy \\
    \textcolor{muted}{Data:} & \datadokumentu \\
    \textcolor{muted}{Wersja aplikacji:} & 0.1.0 (gałąź \code{main}, stan roboczy)
  \end{tabular}
  \vfill
  \begin{tcolorbox}[colback=accentlight,colframe=accentlight,arc=2pt,left=10pt,right=10pt]
    \small
    Dokument składa się z dwóch części. \textbf{Plan testów} opisuje środowisko, narzędzia,
    przypadki testowe pięciu poziomów oraz dane testowe wraz z uzasadnieniem ich wyboru.
    \textbf{Raport z testów} zawiera wyniki wykonania, audyt bezpieczeństwa, analizę
    i zalecenia dla zespołu developerskiego.
  \end{tcolorbox}
\end{titlepage}
\hypersetup{pageanchor=true}

\tableofcontents
\clearpage

% =========================================================================
\part{Plan testów}
% =========================================================================

\section{Wprowadzenie}

\subsection{Cel dokumentu}
Plan określa, co, czym i na jakich danych jest testowane w aplikacji \enquote{Listy zakupów}.
Każdy przypadek ma identyfikator, z którego korzysta część raportowa:
\textbf{TJ} — testy jednostkowe, \textbf{TI} — integracyjne, \textbf{TS} — systemowe,
\textbf{TW} — wydajnościowe, \textbf{TB} — bezpieczeństwa.

\subsection{Przedmiot testów}
Aplikacja webowa dostępna wyłącznie po zalogowaniu. Użytkownik zakłada konto, tworzy listy
zakupów z sugerowaną datą i dodaje pozycje. Listę można udostępnić innej osobie, która może
jedynie zmieniać status pozycji. Frontend: React, Vite, shadcn/ui. Backend: Node.js, Hono,
SQLite z ręcznie pisanymi zapytaniami parametryzowanymi. Zdjęcia pozycji są zapisywane w bazie
jako BLOB.

\begin{table}[H]
\centering\small
\begin{tabularx}{\textwidth}{@{}>{\sffamily\bfseries\color{accent}}l L@{}}
\toprule
\textcolor{black}{ID} & \textsf{\textbf{Wymaganie}}\\
\midrule
R1 & Rejestracja, logowanie i wylogowanie. Bez sesji aplikacja i API są niedostępne.\\
R2 & Pełny CRUD list zakupów. Każda lista ma opcjonalną sugerowaną datę zakupów.\\
R3 & Pozycja ma nazwę wpisaną ręcznie albo wybraną z popularnych produktów (chleb, mleko, woda, ser…), ilość dodatnią i jednostkę z listy \code{szt}, \code{g}, \code{kg}, \code{ml}, \code{l}, \code{opak}.\\
R4 & Pozycja może mieć zdjęcie referencyjne JPEG, PNG lub WebP do 2~MB.\\
R5 & Pozycja kupiona jest na liście przekreślona.\\
R6 & Właściciel udostępnia listę po adresie e-mail. Osoba ze współdzielenia zmienia tylko status, bez dodawania, usuwania i zmiany ilości.\\
R7 & Listy, pozycje, zdjęcia i statusy są trwale zapisane w bazie.\\
R8 & Wymagania niefunkcjonalne: p95 czasu odpowiedzi poniżej 300~ms, mniej niż 1\% błędów, brak podatności o wadze wysokiej.\\
\bottomrule
\end{tabularx}
\caption{Wymagania objęte testami}
\label{tab:wymagania}
\end{table}

\subsection{Metodyka}
Aplikacja powstaje metodyką \textbf{TDD} w cyklu czerwony — zielony — refaktor. Najpierw
powstaje test opisujący jedną regułę i pada, potem najmniejszy kod, który go spełnia, na
końcu porządkowanie przy zielonym teście. Cykle zaczynają się od reguł domeny bez bazy i HTTP,
przechodzą przez API na bazie w pamięci, a kończą na komponentach interfejsu. Test systemowy
w przeglądarce jest testem akceptacyjnym całej ścieżki. Testy wydajności, testy bezpieczeństwa
i skan DAST badają złożoną aplikację i nie są krokiem wewnętrznego cyklu TDD.

\begin{table}[H]
\centering\small
\begin{tabularx}{\textwidth}{@{}l >{\raggedright\arraybackslash}p{3.1cm} L >{\raggedright\arraybackslash}p{4.3cm}@{}}
\toprule
\textsf{\textbf{Poziom}} & \textsf{\textbf{Narzędzie}} & \textsf{\textbf{Zakres}} & \textsf{\textbf{Polecenie}}\\
\midrule
Jednostkowy & Vitest, Testing Library & reguły domeny, komponenty React & \code{pnpm test:unit}\\
Integracyjny & Vitest, \code{app.request()} & trasa HTTP $\rightarrow$ serwis $\rightarrow$ repozytorium $\rightarrow$ SQLite & \code{pnpm test:integration}\\
Systemowy & Playwright, Chromium & cała aplikacja w przeglądarce, dwóch użytkowników & \code{pnpm test:e2e}\\
Wydajnościowy & k6 & logowanie, odczyt list, zmiana statusu pod obciążeniem & \code{pnpm test:perf}\\
Bezpieczeństwa & Vitest, pnpm audit, OWASP ZAP & ataki na API, zależności, skan pasywny & \code{pnpm test:security}\newline\code{pnpm test:zap}\\
\bottomrule
\end{tabularx}
\caption{Poziomy testów i narzędzia}
\end{table}

% -------------------------------------------------------------------------
\section{Środowisko testowe i narzędzia}

\subsection{Sprzęt i system}
\begin{table}[H]
\centering\small
\begin{tabularx}{\textwidth}{@{}l L@{}}
\toprule
\textsf{\textbf{Element}} & \textsf{\textbf{Wartość}}\\
\midrule
System operacyjny & CachyOS Linux, jądro 6.18.50-2-cachyos-lts, x86\_64\\
Procesor & 12 wątków\\
Pamięć RAM & 15 GiB\\
Sieć & pętla zwrotna \code{127.0.0.1}; klient i serwer działają na tej samej maszynie\\
Konteneryzacja & Docker 29.9.0 (tylko dla OWASP ZAP)\\
\bottomrule
\end{tabularx}
\caption{Platforma sprzętowa}
\end{table}

\subsection{Oprogramowanie i narzędzia}
\begin{table}[H]
\centering\small
\begin{tabularx}{\textwidth}{@{}l l L@{}}
\toprule
\textsf{\textbf{Narzędzie}} & \textsf{\textbf{Wersja}} & \textsf{\textbf{Rola}}\\
\midrule
Node.js & 26.10.0 & środowisko uruchomieniowe API, testów i serwera deweloperskiego\\
pnpm & 12.10.1 & menedżer pakietów i workspace, uruchamianie skryptów testowych\\
Hono & 4.13.13 & framework HTTP; \code{app.request()} pozwala testować API bez gniazda sieciowego\\
better-sqlite3 & 12.11.1 & sterownik SQLite; baza \code{:memory:} w testach, plik WAL w pozostałych środowiskach\\
Vitest & 4.1.11 & testy jednostkowe, integracyjne i bezpieczeństwa\\
@vitest/coverage-v8 & 4.1.11 & pomiar pokrycia kodu i egzekwowanie progów\\
Testing Library & 16.3.3 & render komponentów React i symulacja działań użytkownika\\
jsdom & 27.4.0 & przeglądarkowe DOM w Node.js dla testów komponentów\\
Playwright & 1.64.0 & testy systemowe; Chromium build 1248\\
k6 & 2.3.0 & testy wydajnościowe smoke, load, stress\\
OWASP ZAP & 2.17.0 & pasywny skan bezpieczeństwa (obraz \code{zaproxy/zap-stable})\\
pnpm audit & 12.10.1 & audyt znanych podatności w zależnościach\\
TypeScript & 5.9 & kontrola typów jako dodatkowa bramka jakości\\
\bottomrule
\end{tabularx}
\caption{Oprogramowanie użyte do testów}
\end{table}

\subsection{Konfiguracja środowisk}
Każdy poziom testów ma osobną bazę danych, dzięki czemu przebiegi nie wpływają na siebie
i na dane deweloperskie.

\begin{table}[H]
\centering\small
\begin{tabularx}{\textwidth}{@{}l l l L@{}}
\toprule
\textsf{\textbf{Środowisko}} & \textsf{\textbf{Adres}} & \textsf{\textbf{Baza}} & \textsf{\textbf{Ustawienia}}\\
\midrule
Jednostkowe i integracyjne & brak gniazda & \code{:memory:} & nowa baza na każdy test; limit logowań 1000, ciasteczka bez \code{Secure}\\
Systemowe & \code{:3000} API, \code{:5173} front & \code{e2e.sqlite} & plik usuwany przed przebiegiem; serwery stawia Playwright\\
Wydajnościowe & \code{:3099} & \code{perf.sqlite} & plik tworzony od zera; \code{LOGIN\_RATE\_LIMIT\_MAX=100000}\\
Skan ZAP & \code{:3000} & \code{app.sqlite} & kontener w sieci hosta, bo API nasłuchuje tylko na \code{127.0.0.1}\\
\bottomrule
\end{tabularx}
\caption{Konfiguracja środowisk testowych}
\end{table}

% -------------------------------------------------------------------------
\section{Kryteria wejścia i wyjścia}

\textbf{Kryteria wejścia:} zależności zainstalowane (\code{pnpm install}), kod przechodzi
kontrolę typów, dla testów systemowych zainstalowana przeglądarka Chromium, dla testów
wydajności dostępny k6, dla skanu ZAP działający Docker.

\textbf{Kryteria wyjścia:}
\begin{enumerate}
  \item wszystkie przypadki TJ, TI, TS i TB zakończone wynikiem pozytywnym;
  \item pokrycie kodu nie niższe niż progi zapisane w konfiguracji Vitest:
        API 90\% instrukcji, 85\% gałęzi, 90\% funkcji, 95\% linii;
        front 90\% instrukcji, 80\% gałęzi, 90\% funkcji, 90\% linii;
  \item w scenariuszach smoke i load: p95 poniżej 300~ms i odsetek błędów poniżej 1\%.
        Scenariusz stress ma ten próg przekroczyć, przy błędach nadal poniżej 1\%,
        żeby wskazać granicę obciążenia;
  \item brak podatności o wadze \emph{high} lub \emph{critical} w \code{pnpm audit};
  \item brak alertów typu FAIL w skanie OWASP ZAP baseline.
\end{enumerate}

% -------------------------------------------------------------------------
\section{Testy jednostkowe — przypadki testowe}

\subsection{Zakres i techniki}
Testy jednostkowe sprawdzają czystą logikę bez bazy i bez sieci oraz komponenty interfejsu
w izolacji od backendu (żądania \code{fetch} są zastępowane atrapami). Do projektowania danych
użyto \textbf{klas równoważności} (poprawna i niepoprawna jednostka), \textbf{wartości
brzegowych} (ilość 0 i 1\,000\,001, rozmiar zdjęcia 2~MB + 1~bajt, hasło 7 znaków) oraz
\textbf{tablicy decyzyjnej} dla uprawnień (tabela~\ref{tab:uprawnienia}).

\begin{table}[H]
\centering\small
\begin{tabular}{@{}lccc@{}}
\toprule
\textsf{\textbf{Działanie}} & \textsf{\textbf{Właściciel}} & \textsf{\textbf{Osoba ze współdzielenia}} & \textsf{\textbf{Obca osoba}}\\
\midrule
Podgląd listy & tak & tak & nie (404)\\
Zmiana statusu pozycji & tak & tak & nie (404)\\
Dodanie, edycja, usunięcie pozycji & tak & nie (403) & nie (404)\\
Edycja i usunięcie listy, udostępnianie & tak & nie (403) & nie (404)\\
\bottomrule
\end{tabular}
\caption{Tablica decyzyjna uprawnień (R6)}
\label{tab:uprawnienia}
\end{table}

\subsection{Backend — reguły domeny}
\begin{przypadki}
\plik{apps/api/src/domain/units.test.ts}
TJ-01 & Jednostka z listy (R3) & każda z \code{szt}, \code{g}, \code{kg}, \code{ml}, \code{l}, \code{opak}; \code{tona}; pusty napis & jednostki z listy przyjęte, \code{tona} i pusty napis odrzucone\\
TJ-02 & Wartości brzegowe ilości (R3) & 1; 1,5; 0; $-2$; NaN; $+\infty$; 1\,000\,001 & 1 i 1,5 przyjęte, pozostałe odrzucone\\
\plik{apps/api/src/domain/schemas.test.ts}
TJ-03 & Walidacja rejestracji (R1) & hasło \code{haslo1234} (9 znaków); hasło \code{krotkie} (7 znaków) & pierwsze przyjęte, drugie odrzucone\\
TJ-04 & Data listy (R2) & \code{2026-02-31}; \code{2026-10-12}; pusta zmiana listy & nieistniejąca data i pusta zmiana odrzucone, poprawna data przyjęta\\
TJ-05 & Schemat pozycji (R3) & ilość \code{"2"} jako tekst z formularza; ilość 0; jednostka \code{tona} & tekst zamieniony na liczbę i przyjęty, pozostałe odrzucone\\
\plik{apps/api/src/domain/items.test.ts}
TJ-06 & Kolejność pozycji (R3) & trzy pozycje: pozycja 2; pozycja 1 z lutego; pozycja 1 ze stycznia & kolejność \code{a}, \code{c}, \code{b}: po pozycji, przy remisie po dacie\\
TJ-07 & Podsumowanie listy (R5) & statusy: kupione, niekupione, kupione & razem 3, kupione 2, pozostało 1\\
\plik{apps/api/src/domain/rateLimit.test.ts}
TJ-08 & Okno limitu prób (R8) & limit 2, okno 500~ms; próby w chwilach 1000, 1100, 1200, 1600~ms & dwie pierwsze dozwolone, trzecia zablokowana, czwarta po resecie okna dozwolona\\
\plik{apps/api/src/domain/image.test.ts}
TJ-09 & Poprawny PNG (R4) & plik PNG 1$\times$1 piksel, 70~B & przyjęty\\
TJ-10 & Niepoprawne zdjęcie (R4) & typ \code{text/plain}; bajty \code{01 02 03 04} jako PNG; pusty plik; PNG o rozmiarze 2~MB + 1~B & wszystkie odrzucone z komunikatem\\
TJ-11 & Podpisy JPEG i WebP (R4) & \code{FF D8 FF}; \code{00 00 00} jako JPEG; nagłówek \code{RIFF…WEBP}; ucięty WebP & poprawne podpisy przyjęte, błędne odrzucone\\
\plik{apps/api/src/domain/permissions.test.ts}
TJ-12 & Prawa właściciela (R6) & właściciel listy & może edytować pozycje i zmieniać status\\
TJ-13 & Prawa osoby ze współdzielenia (R6) & użytkownik z listy udostępnień & zmienia tylko status, edycja zabroniona\\
TJ-14 & Obca osoba (R6) & użytkownik spoza listy & brak podglądu, edycji i zmiany statusu\\
\end{przypadki}

\subsection{Frontend — komponenty i logika klienta}
\begin{przypadki}
\plik{apps/web/src/components/ItemRow.test.tsx}
TJ-15 & Przekreślenie kupionej pozycji (R5) & pozycja \code{isPurchased: true} & etykieta ma klasę \code{line-through}\\
TJ-16 & Widok osoby ze współdzielenia (R6) & \code{canEdit: false} & brak przycisków edycji i usuwania\\
TJ-17 & Zgłoszenie zmiany statusu (R5) & kliknięcie pola wyboru & \code{onToggle(true)}\\
\plik{apps/web/src/components/AddItemForm.test.tsx}
TJ-18 & Dodanie pozycji z domyślną jednostką (R3) & nazwa \code{mleko}, ilość 2 & wysłane \code{\{mleko, 2, szt, bez zdjęcia\}}\\
TJ-19 & Wybór produktu z podpowiedzi i zdjęcie (R3, R4) & wpis \code{mle}, wybór \code{mleko}, plik PNG & wysłane \code{mleko}, ilość 1 i plik\\
TJ-20 & Nazwa spoza katalogu (R3) & wpis \code{tofu}, opcja \enquote{Dodaj „tofu”} & wysłana nazwa \code{tofu}\\
TJ-21 & Pusta nazwa (R3) & wysłanie pustego formularza & brak wysyłki, widoczny komunikat błędu\\
\plik{apps/web/src/components/ProtectedRoute.test.tsx}
TJ-22 & Strona chroniona bez sesji (R1) & brak użytkownika & przekierowanie na logowanie\\
TJ-23 & Strona chroniona z sesją (R1) & zalogowany Jan & widoczna treść list\\
TJ-24 & Sprawdzanie sesji (R1) & \code{loading: true} & wskaźnik ładowania\\
TJ-25 & Logowanie dla gościa (R1) & brak użytkownika na \code{/login} & widoczny ekran logowania\\
TJ-26 & Logowanie przy aktywnej sesji (R1) & zalogowany Jan na \code{/login} & przekierowanie na listy\\
\plik{apps/web/src/components/DateField.test.tsx}
TJ-27 & Formatowanie daty (R2) & \code{null}; \code{wtorek}; \code{2026-10-16} & \enquote{Bez daty}; tekst bez zmian; data po polsku\\
TJ-28 & Wybór daty polem i kalendarzem (R2) & wpis \code{2026-10-20}; kliknięcie dnia w kalendarzu & \code{onChange} wywołane w obu przypadkach\\
\plik{apps/web/src/components/ShareDialog.test.tsx}
TJ-29 & Udostępnienie i odebranie dostępu (R6) & e-mail \code{ania@example.com}, przycisk \enquote{Odbierz} & \code{onShare} z adresem, \code{onUnshare('2')}\\
TJ-30 & Nieudane udostępnienie (R6) & \code{onShare} rzuca błąd & adres zostaje w polu\\
\plik{apps/web/src/components/AppHeader.test.tsx}
TJ-31 & Wylogowanie (R1) & kliknięcie \enquote{Wyloguj} & wywołane \code{logout}, przejście na logowanie\\
\plik{apps/web/src/auth/AuthProvider.test.tsx}
TJ-32 & Cykl sesji w kliencie (R1) & odpowiedzi \code{me}, \code{login}, \code{register}, \code{logout} & nazwa użytkownika zmienia się po każdym kroku, na końcu gość\\
TJ-33 & Brak sesji (R1) & \code{/api/auth/me} zwraca 401 & stan gościa\\
TJ-34 & Użycie poza providerem & render bez \code{AuthProvider} & błąd \enquote{useAuth poza AuthProvider}\\
\plik{apps/web/src/lib/api.test.ts, apps/web/src/lib/schemas.test.ts}
TJ-35 & Nagłówek JSON & ciało tekstowe & \code{content-type: application/json}, zwrócone ciało\\
TJ-36 & Formularz i pusta odpowiedź & \code{FormData}, odpowiedź 204 & brak nadpisanego typu, wynik \code{undefined}\\
TJ-37 & Obsługa błędów API & 400 z \code{error}; 500 bez JSON & komunikat z serwera; komunikat ogólny\\
TJ-38 & Schemat pozycji w kliencie (R3) & jednostka \code{tona} & odrzucona\\
\plik{apps/web/src/pages/auth-pages.test.tsx}
TJ-39 & Udane logowanie (R1) & \code{jan@example.com} / \code{haslo1234} & wywołane \code{login}, przejście na listy\\
TJ-40 & Błędne hasło (R1) & \code{login} rzuca \code{ApiError(401)} & widoczny komunikat z serwera\\
TJ-41 & Pusty formularz logowania (R1) & brak danych & brak wywołania \code{login}\\
TJ-42 & Udana rejestracja (R1) & Jan, e-mail, hasło 9 znaków & wywołane \code{register}, przejście na listy\\
TJ-43 & Błąd sieci przy rejestracji (R1) & zwykły \code{Error} & komunikat \enquote{Nie udało się utworzyć konta}\\
\plik{apps/web/src/pages/lists-flow.test.tsx, apps/web/src/App.test.tsx}
TJ-44 & Listy własne i udostępnione (R2, R6) & listy \code{Piątek} i \code{Od Ani}, nowa \code{Sobota} & obie sekcje widoczne; po utworzeniu widok nowej, pustej listy\\
TJ-45 & Pusty stan (R2) & brak list & komunikaty pustych sekcji\\
TJ-46 & Pełna obsługa listy przez właściciela (R2–R6) & odhaczenie, dodanie \code{chleb}, udostępnienie, zmiana nazwy, edycja i usunięcie pozycji, usunięcie listy & każde działanie wysyła żądanie, na końcu powrót do list\\
TJ-47 & Błąd wczytania listy & odpowiedź 500 & komunikat \enquote{Nie udało się otworzyć listy.}\\
TJ-48 & Widok osoby ze współdzielenia (R6) & rola \code{shared}, zmiana statusu kończy się 500 & notka o prawach, brak edycji, interfejs wraca do stanu sprzed zmiany\\
TJ-49 & Nieznana ścieżka bez sesji (R1) & adres \code{/nie-ma}, \code{me} zwraca 401 & ekran logowania\\
\end{przypadki}

% -------------------------------------------------------------------------
\section{Testy integracyjne — przypadki testowe}

Testy integracyjne wysyłają żądania HTTP do aplikacji Hono przez \code{app.request()}.
Żądanie przechodzi przez middleware (nagłówki, limit rozmiaru, sesja), trasę, serwis
i repozytorium do prawdziwej bazy SQLite w pamięci z pełnym schematem. Repozytoria nie mają
osobnych testów: ich zachowanie wymuszają poniższe przypadki.

\begin{przypadki}
\plik{apps/api/src/integration/auth.test.ts}
TI-01 & Rejestracja, sesja i wylogowanie (R1) & Jan, \code{jan@example.com}, \code{haslo1234} & 201 i ciasteczko \code{session}; \code{/me} 200; wylogowanie 204; \code{/me} 401\\
TI-02 & Brak wycieku istnienia konta (R1) & nieznany e-mail; znany e-mail ze złym hasłem & oba 401 z identyczną treścią\\
TI-03 & Unikalny e-mail bez względu na wielkość liter (R1) & druga rejestracja \code{Jan@example.com} & 409\\
TI-04 & Wymóg sesji (R1) & \code{GET /api/lists} bez ciasteczka & 401\\
\plik{apps/api/src/integration/lists.test.ts}
TI-05 & Pełny CRUD listy i pozycji ze zdjęciem (R2–R5, R7) & lista \code{Zakupy na piątek} 2026-10-16; pozycja \code{mleko} 2~l z PNG w \code{multipart} & 201; zmiana nazwy; zdjęcie zwrócone bajt w bajt z typem \code{image/png}; edycja; status kupione; usunięcie 204; lista 404\\
TI-06 & Kaskadowe usuwanie (R7) & lista z pozycją \code{chleb} & po usunięciu listy w tabeli \code{items} 0 wierszy\\
TI-07 & Katalog i jednostki (R3) & \code{/popular-products} bez sesji i z sesją; \code{/units} & 401; lista zawiera \code{chleb}; jednostki dokładnie \code{szt, g, kg, ml, l, opak}\\
\plik{apps/api/src/integration/sharing.test.ts}
TI-08 & Prawa osoby ze współdzielenia (R6) & Jan udostępnia listę z pozycją \code{woda} 6~szt Ani & Ania: dodanie, edycja, usunięcie 403; zmiana statusu 200; Jan widzi status kupione i licznik 1\\
TI-09 & Obca lista i udostępnienie sobie (R6) & Obcy otwiera listę Jana; Jan udostępnia sam sobie & 404; 400\\
\plik{apps/api/src/integration/edges.test.ts}
TI-10 & Sprawdzenie stanu i nieznana ścieżka & \code{/api/health}; \code{/api/brak} & \code{\{ok: true\}}; 404\\
TI-11 & Logowanie i ciasteczko \code{Secure} (R1) & logowanie istniejącego konta przy \code{secureCookies} & 200 i flaga \code{Secure}; wylogowanie bez ciasteczka 204\\
TI-12 & Złe ciało żądania i limit przez proxy (R1, R8) & ciało \code{null}; JSON \code{\{}; limit 1 i \code{X-Forwarded-For: 10.0.0.8} & 400; 400; druga próba z tego adresu 429\\
TI-13 & Za duże żądanie i awaria bazy & ciało 3~MB + 1~B; zamknięta baza & 413; 500 z ogólnym komunikatem\\
TI-14 & Zarządzanie udostępnieniami (R2, R6) & usunięcie daty; udostępnienie nieznanemu; ponowne udostępnienie; odebranie przez Anię i przez Jana & data \code{null}; 404; 409; 403; 204, ponownie 404\\
TI-15 & Wymiana zdjęcia (R4) & brak zdjęcia; pusty formularz; PNG; usunięcie przez Anię i Jana; zdjęcie 2~MB + 1~B & 404; 400; \code{hasImage: true}; 403; 204; 400\\
TI-16 & Plikowa baza i seed (R7) & nowy katalog tymczasowy, dwukrotny seed & tryb WAL; produkty zasiane, drugi seed nie dubluje danych\\
\end{przypadki}

% -------------------------------------------------------------------------
\section{Dane do testów integracyjnych}

\subsection{Zestaw danych}
{\small
\begin{xltabular}{\textwidth}{@{}l L L@{}}
\toprule
\textsf{\textbf{Kategoria}} & \textsf{\textbf{Dane}} & \textsf{\textbf{Uzasadnienie}}\\
\midrule
\endfirsthead
\toprule
\textsf{\textbf{Kategoria}} & \textsf{\textbf{Dane}} & \textsf{\textbf{Uzasadnienie}}\\
\midrule
\endhead
\bottomrule
\endfoot
Baza & SQLite \code{:memory:} ze schematem produkcyjnym i 20 popularnymi produktami & każdy test zaczyna od czystego stanu; ten sam schemat i te same ograniczenia \code{CHECK} co w produkcji\\
Użytkownicy & Jan \code{jan@example.com} (właściciel), Ania \code{ania@example.com} (współdzielenie), Obcy \code{obcy@example.com} & po jednym przedstawicielu każdej roli z tablicy decyzyjnej uprawnień\\
Hasło & \code{haslo1234} (9 znaków) & minimalnie powyżej granicy 8 znaków; poprawne dla walidatora\\
Wielkość liter & \code{Jan@example.com} & sprawdza normalizację e-maila przy unikalności konta\\
Listy & \code{Zakupy na piątek}, data 2026-10-16; \code{Na dziś} bez daty & obie klasy: lista z datą i bez daty\\
Pozycje & \code{mleko} 2~l, \code{woda} 6~szt, \code{ser} 1~opak, \code{chleb} 1~szt & produkty z katalogu i różne jednostki, w tym niedomyślne\\
Zdjęcie & PNG 1$\times$1 piksel, 70~B & najmniejszy poprawny plik; szybkie testy i porównanie bajt w bajt\\
Wartości brzegowe & zdjęcie 2~MB + 1~B; ciało 3~MB + 1~B & pierwszy bajt ponad limit zdjęcia i limit ciała żądania\\
Złe dane & ciało \code{null}, JSON \code{\{}, nieznany e-mail, nieznane ID listy & ścieżki błędów, które klient może wywołać\\
Proxy & \code{X-Forwarded-For: 10.0.0.8, 10.0.0.1}, limit logowań 1 & pokazuje, od czego zależy klucz limitu\\
\end{xltabular}
\vspace{-1.2em}\captionof{table}{Dane do testów integracyjnych}}

\subsection{Dlaczego taki zestaw}
\begin{itemize}
  \item \textbf{Izolacja i powtarzalność.} Baza w pamięci powstaje dla każdego testu, więc
        kolejność wykonania nie ma znaczenia, a cały zestaw trwa około 3~s.
  \item \textbf{Prawdziwy silnik zamiast atrapy.} Zapytania SQL, klucze obce
        i \code{ON DELETE CASCADE} wykonuje SQLite, więc test TI-06 sprawdza rzeczywiste
        zachowanie bazy, a nie założenie o nim.
  \item \textbf{Pokrycie ról.} Trzy konta odpowiadają trzem kolumnom tablicy decyzyjnej;
        mniej kont nie pozwoliłoby odróżnić odpowiedzi 403 od 404.
  \item \textbf{Wartości brzegowe zamiast losowych.} Pliki i ciała żądań są o jeden bajt
        większe od limitu, bo tam najczęściej powstają błędy typu \enquote{off by one}.
  \item \textbf{Czytelność.} Polskie nazwy produktów z katalogu odpowiadają realnemu użyciu
        aplikacji i ułatwiają analizę nieudanego testu.
\end{itemize}

% -------------------------------------------------------------------------
\section{Testy systemowe — przypadki testowe}

Testy systemowe uruchamiają rzeczywiste API i serwer frontendu, a następnie sterują
przeglądarką Chromium jak użytkownik. Scenariusz TS-02 używa dwóch niezależnych kontekstów
przeglądarki, czyli dwóch osobnych sesji.

\begin{przypadki}
TS-01 & Logowanie istniejącym kontem (R1) & Ola, \code{login.<znacznik czasu>@example.com}, \code{haslo1234} & po rejestracji i wylogowaniu ponowne logowanie otwiera stronę list\\
TS-02 & Właściciel zarządza listą, osoba ze współdzielenia zmienia tylko status (R1–R7) & Jan i Ania z unikalnymi adresami; lista \code{Zakupy na piątek}, 2026-10-16; \code{mleko} 2~l z \code{pixel.png} & kroki i oczekiwane wyniki w tabeli~\ref{tab:ts02}\\
TS-03 & Weryfikacja ręczna w przeglądarce (R1–R5, R7) & konto \code{laboratorium@example.com}; lista \code{Zakupy na piątek} wybrana kalendarzem; \code{mleko} 2~l & rejestracja, lista z datą, przekreślenie, stan 1/1 po odświeżeniu; \code{/lists} po wylogowaniu przekierowuje na logowanie\\
\end{przypadki}

\begin{table}[H]
\centering\footnotesize
\begin{tabularx}{\textwidth}{@{}r l L L@{}}
\toprule
\textsf{\textbf{Krok}} & \textsf{\textbf{Kto}} & \textsf{\textbf{Działanie}} & \textsf{\textbf{Oczekiwany wynik}}\\
\midrule
1 & Ania, Jan & rejestracja w osobnych kontekstach & obie osoby widzą stronę list\\
2 & Jan & nowa lista, otwarcie kalendarza, data 2026-10-16 & kalendarz widoczny; tytuł \code{Zakupy na piątek}\\
3 & Jan & pozycja \code{mleko}, ilość 2, jednostka \code{l}, zdjęcie & pozycja z miniaturą zdjęcia\\
4 & Jan & odhaczenie pozycji & etykieta przekreślona (\code{line-through})\\
5 & Jan & zmiana nazwy na \code{Zakupy na sobotę} & nowy tytuł\\
6 & Jan & udostępnienie Ani po e-mailu & adres Ani na liście udostępnień\\
7 & Ania & otwarcie udostępnionej listy & notka o prawach; brak formularza dodawania i przycisku usuwania\\
8 & Ania & zmiana statusu pozycji & przekreślenie znika\\
9 & Jan & odświeżenie strony & widzi zmianę Ani: pozycja nieprzekreślona\\
10 & Jan & usunięcie pozycji z potwierdzeniem & lista bez pozycji\\
11 & Jan & usunięcie listy z potwierdzeniem & powrót do list; listy \code{Zakupy na sobotę} brak\\
\bottomrule
\end{tabularx}
\caption{Kroki scenariusza TS-02}
\label{tab:ts02}
\end{table}

% -------------------------------------------------------------------------
\section{Dane do testów systemowych}

{\small
\begin{xltabular}{\textwidth}{@{}l L L@{}}
\toprule
\textsf{\textbf{Kategoria}} & \textsf{\textbf{Dane}} & \textsf{\textbf{Uzasadnienie}}\\
\midrule
\endfirsthead
\toprule
\textsf{\textbf{Kategoria}} & \textsf{\textbf{Dane}} & \textsf{\textbf{Uzasadnienie}}\\
\midrule
\endhead
\bottomrule
\endfoot
Baza & \code{e2e.sqlite} usuwany przez \code{global-setup.ts} przed przebiegiem & start z czystą bazą, bez danych z poprzednich uruchomień\\
Konta & \code{jan.<czas>@}, \code{ania.<czas>@}, \code{login.<czas>@example.com} & unikalne adresy pozwalają uruchamiać testy wielokrotnie bez konfliktu 409\\
Hasło & \code{haslo1234} & poprawne według reguły 8 znaków, takie samo jak w testach integracyjnych\\
Lista & \code{Zakupy na piątek} $\rightarrow$ \code{Zakupy na sobotę}, data 2026-10-16 & polskie znaki diakrytyczne sprawdzają kodowanie w całym łańcuchu; zmiana nazwy sprawdza edycję\\
Pozycja & \code{mleko}, 2, \code{l} & produkt z katalogu podpowiedzi; jednostka inna niż domyślna \code{szt} wymusza użycie listy rozwijanej\\
Zdjęcie & \code{fixtures/pixel.png}, 70~B & ścieżka przesłania pliku przez formularz bez spowalniania testu\\
Dwie sesje & dwa konteksty Playwright & jedyny sposób, by w jednym teście zobaczyć prawa obu ról i propagację zmian\\
\end{xltabular}
\vspace{-1.2em}\captionof{table}{Dane do testów systemowych}}

Dane odwzorowują typowe użycie aplikacji: jedna osoba planuje zakupy, druga odhacza
produkty w sklepie. Wybrano najkrótszą ścieżkę, która przechodzi przez każde wymaganie
funkcjonalne R1–R7. Przypadki brzegowe i błędne zostały sprawdzone niżej w piramidzie,
gdzie są szybsze i stabilniejsze.

% -------------------------------------------------------------------------
\section{Testy wydajnościowe — przypadki testowe}

\subsection{Model obciążenia}
Każdy scenariusz k6 loguje się raz w fazie \code{setup}, a następnie w pętli wykonuje
\code{GET /api/lists} oraz \code{PATCH /api/items/:id/status}. Osobna grupa wirtualnych
użytkowników (VU) loguje się cyklicznie, bo logowanie liczy skrót bcrypt i jest najdroższą
operacją API. Skrypt \code{tests/perf/run.sh} uruchamia API na porcie 3099 na nowej bazie,
seeduje dane, wykonuje trzy scenariusze i zatrzymuje serwer.

\begin{przypadki}
TW-01 & Smoke — czy system działa pod minimalnym ruchem (R8) & 1 VU przez 10~s; przerwa 0,2~s & p95 < 300~ms; błędy < 1\%; wszystkie sprawdzenia 200\\
TW-02 & Load — typowe obciążenie (R8) & 8 VU przegląda i odhacza, 2 VU loguje się; 20~s & jak wyżej\\
TW-03 & Stress — rosnące obciążenie, aż do przekroczenia progu (R8) & przegląd: 0 $\rightarrow$ 20 $\rightarrow$ 35 $\rightarrow$ 40 $\rightarrow$ 0 VU; logowanie: 0 $\rightarrow$ 6 $\rightarrow$ 12 $\rightarrow$ 16 $\rightarrow$ 0 VU; 37~s & p95 powyżej 300~ms; błędy nadal poniżej 1\%; sprawdzenia 200\\
\end{przypadki}

\subsection{Dane}
Seed \code{tests/perf/seed.mjs} tworzy konto \code{perf@example.com} oraz 5 list po 10 pozycji,
razem 50 pozycji. Odpowiedź \code{GET /api/lists} zawiera więc liczniki wyliczane zapytaniem
agregującym na realnej liczbie wierszy. Limit logowań w tym środowisku podniesiono do 100\,000
na minutę. Przy domyślnych 10 próbach na minutę scenariusze mierzyłyby odpowiedzi 429
zamiast wydajności.

% -------------------------------------------------------------------------
\section{Audyt bezpieczeństwa — plan}

\subsection{Zakres i metody}
Audyt obejmuje API i zależności projektu. Wykorzystuje cztery metody:
\begin{enumerate}
  \item \textbf{testy automatyczne ataków} (Vitest, przypadki TB-01–TB-07);
  \item \textbf{audyt zależności} \code{pnpm audit --audit-level=high};
  \item \textbf{skan pasywny DAST} OWASP ZAP baseline na działającym API;
  \item \textbf{przegląd kodu} pod kątem OWASP Top 10 (2021).
\end{enumerate}

\begin{przypadki}
TB-01 & Nagłówki ochronne i flagi ciasteczka & \code{/api/health}; rejestracja przy \code{secureCookies} & \code{X-Content-Type-Options: nosniff}, \code{X-Frame-Options: SAMEORIGIN}; ciasteczko \code{HttpOnly}, \code{SameSite=Lax}, \code{Secure}\\
TB-02 & Dostęp bez sesji & \code{/api/lists}, \code{/api/auth/me}, \code{/api/popular-products} & 401 dla każdego\\
TB-03 & Zgadywanie hasła & 4 próby ze złym hasłem przy limicie 3 & 3 razy 401, czwarta 429\\
TB-04 & Wstrzyknięcie SQL & nazwa listy \code{'; DROP TABLE lists; --}; e-mail \code{jan@example.com' OR '1'='1} & nazwa zapisana jako tekst; logowanie 400; tabela i konto nienaruszone\\
TB-05 & XSS zapisany w danych & nazwa pozycji \code{<img src=x onerror=alert(1)>} & odpowiedź JSON, wartość bez zmian, brak dokumentu HTML\\
TB-06 & IDOR — cudze zasoby & Obcy pyta o listę, zdjęcie i status pozycji Jana & 404 dla każdego; właściciel dalej ma dostęp\\
TB-07 & Złośliwe przesłanie pliku & \code{text/html}; tekst udający PNG; PNG 2~MB + 1~B & 400 dla każdego; komunikat o limicie 2~MB\\
\end{przypadki}

\subsection{Przegląd kodu — lista kontrolna}
\begin{table}[H]
\centering\small
\begin{tabularx}{\textwidth}{@{}l L l@{}}
\toprule
\textsf{\textbf{Kategoria OWASP}} & \textsf{\textbf{Co jest sprawdzane}} & \textsf{\textbf{Metoda}}\\
\midrule
A01 Broken Access Control & kontrola roli przy każdej operacji, 404 dla obcych & TB-06, TI-08, TI-09, przegląd\\
A02 Cryptographic Failures & bcrypt, losowość tokenu sesji, przechowywanie tokenu & przegląd\\
A03 Injection & zapytania parametryzowane, odpowiedzi JSON & TB-04, TB-05, przegląd\\
A04 Insecure Design & limit prób logowania, jego klucz i zasięg & TB-03, TI-12, przegląd\\
A05 Security Misconfiguration & nagłówki, cache, ujawnianie błędów & TB-01, TI-13, ZAP\\
A06 Vulnerable Components & znane podatności zależności & \code{pnpm audit}\\
A07 Identification and Authentication Failures & sesje, wylogowanie, komunikaty logowania & TB-02, TI-01, TI-02\\
A08 Software and Data Integrity & weryfikacja treści plików, polityka łańcucha dostaw & TB-07, przegląd\\
A09 Logging and Monitoring & rejestrowanie zdarzeń bezpieczeństwa & przegląd\\
A10 SSRF & wychodzące żądania serwera & przegląd (brak takich żądań)\\
\bottomrule
\end{tabularx}
\caption{Lista kontrolna przeglądu kodu}
\end{table}

% =========================================================================
\part{Raport z testów}
% =========================================================================

\section{Podsumowanie wykonania}

Wszystkie testy wykonano 12 października 2026 w środowisku z rozdziału~2.

\begin{table}[H]
\centering\small
\begin{tabularx}{\textwidth}{@{}l r r r r L@{}}
\toprule
\textsf{\textbf{Poziom}} & \textsf{\textbf{Przypadki}} & \textsf{\textbf{Zal.}} & \textsf{\textbf{Niezal.}} & \textsf{\textbf{Czas}} & \textsf{\textbf{Wynik}}\\
\midrule
Jednostkowe — backend & 14 & 14 & 0 & \multirow{2}{*}{5,8 s\textsuperscript{a}} & \zal\\
Integracyjne & 16 & 16 & 0 & & \zal\\
Jednostkowe — frontend & 35 & 35 & 0 & 6,1 s & \zal\\
Systemowe (automatyczne) & 2 & 2 & 0 & 4,7 s & \zal\\
Systemowe (ręczne) & 1 & 1 & 0 & — & \zal\ po poprawce D5\\
Wydajnościowe & 3 & 2 & 1 & 37 s & smoke i load \zal; stress ponad próg\\
Bezpieczeństwa (Vitest) & 7 & 7 & 0 & 1,0 s & \zal\\
\code{pnpm audit} & 1 & 1 & 0 & < 1 s & \zal: 0 podatności\\
OWASP ZAP baseline & 67 reguł & 66 & 0 FAIL & ok. 2 min & \zal: 1 ostrzeżenie\\
\midrule
\textbf{Razem testy} & \textbf{78} & \textbf{77} & \textbf{1} & & stress ponad próg p95\\
\bottomrule
\end{tabularx}
\par\smallskip
\raggedright\footnotesize\textsuperscript{a} Testy jednostkowe backendu i integracyjne wykonują się razem w przebiegu z pomiarem pokrycia.
\caption{Zbiorcze wyniki}
\end{table}

\begin{werdykt}[Werdykt]
Wymagania funkcjonalne R1–R7 są potwierdzone na trzech poziomach testów. Wymaganie R8
jest spełnione w smoke i load, do 10 wirtualnych użytkowników. Stress przy 56 użytkownikach
(16 logujących się i 40 przeglądających) przekracza próg: p95 wynosi 653~ms, a błędów HTTP
nie ma. Przegląd kodu wykazał dwa ustalenia o wadze średniej (B1, B2), które należy usunąć
przed wdrożeniem produkcyjnym.
\end{werdykt}

% -------------------------------------------------------------------------
\section{Wyniki testów jednostkowych i integracyjnych}

\subsection{Wykonanie}
Wszystkie 65 przypadków TJ-01–TJ-49 i TI-01–TI-16 zakończyło się wynikiem pozytywnym
(30 testów w 10 plikach dla API, 35 testów w 12 plikach dla frontu).

\subsection{Pokrycie kodu}
Pokrycie mierzono narzędziem V8. Progi są zapisane w konfiguracji Vitest obu pakietów,
więc spadek poniżej nich kończy \code{pnpm test:unit} błędem. Poza pomiarem pozostają tylko
pliki startowe procesu: \code{index.ts} w API i \code{main.tsx} we froncie.

\begin{table}[H]
\centering\small
\begin{tabular}{@{}l rrrr rrrr@{}}
\toprule
& \multicolumn{4}{c}{\textsf{\textbf{API}}} & \multicolumn{4}{c}{\textsf{\textbf{Front}}}\\
\cmidrule(lr){2-5}\cmidrule(l){6-9}
& instr. & gał. & funk. & linie & instr. & gał. & funk. & linie\\
\midrule
Przed uzupełnieniem testów\textsuperscript{b} & 80,0 & 62,9 & 81,3 & 84,0 & 18,9 & 20,3 & 18,8 & 19,1\\
\textbf{Po uzupełnieniu} & \textbf{95,8} & \textbf{89,1} & \textbf{95,8} & \textbf{99,1} & \textbf{97,0} & \textbf{82,2} & \textbf{96,1} & \textbf{97,7}\\
Próg & 90 & 85 & 90 & 95 & 90 & 80 & 90 & 90\\
\bottomrule
\end{tabular}
\par\smallskip
\raggedright\footnotesize\textsuperscript{b} Pierwszy pomiar obejmował również pliki startowe, więc wartości są nieco zaniżone.
\caption{Pokrycie kodu w procentach}
\end{table}

\begin{center}
\begin{tikzpicture}
\begin{axis}[
  ybar=3pt, bar width=14pt, width=0.95\textwidth, height=5.4cm,
  ymin=0, ymax=112, ytick={0,25,50,75,100}, ylabel={\% pokrycia},
  symbolic x coords={API instr.,API gał.,Front instr.,Front gał.},
  xtick=data, enlarge x limits=0.18,
  legend style={at={(0.5,1.04)},anchor=south,legend columns=2,draw=none},
  tick label style={font=\small}, label style={font=\small},
  nodes near coords, every node near coord/.append style={font=\scriptsize},
  nodes near coords style={/pgf/number format/.cd,fixed,precision=1,use comma},
  axis line style={draw=muted}, major grid style={draw=muted!25}, ymajorgrids]
\addplot[fill=muted!40,draw=muted] coordinates {(API instr.,80.0) (API gał.,62.9) (Front instr.,18.9) (Front gał.,20.3)};
\addplot[fill=accent,draw=accent] coordinates {(API instr.,95.8) (API gał.,89.1) (Front instr.,97.0) (Front gał.,82.2)};
\legend{przed, po}
\end{axis}
\end{tikzpicture}
\end{center}

\subsection{Obszary o słabszym pokryciu}
\begin{itemize}
  \item \code{AuthProvider.tsx}, 62,5\% gałęzi: nie jest sprawdzone odmontowanie komponentu
        przed odpowiedzią \code{/api/auth/me}.
  \item \code{ListDetailPage.tsx}, 75,6\% gałęzi: w oknie edycji pozycji nie są zmieniane
        pola ilości i jednostki, wysyłane są tylko istniejące wartości.
  \item \code{apps/api/src/http.ts}, 60\% gałęzi: brak przypadku, w którym błąd walidacji
        nie ma komunikatu.
\end{itemize}

% -------------------------------------------------------------------------
\section{Wyniki testów systemowych}

\begin{table}[H]
\centering\small
\begin{tabularx}{\textwidth}{@{}l L r l@{}}
\toprule
\textsf{\textbf{ID}} & \textsf{\textbf{Przypadek}} & \textsf{\textbf{Czas}} & \textsf{\textbf{Wynik}}\\
\midrule
TS-01 & Logowanie istniejącym kontem & 0,8 s & \zal\\
TS-02 & Właściciel zarządza listą, osoba ze współdzielenia zmienia tylko status & 2,6 s & \zal\\
TS-03 & Weryfikacja ręczna w przeglądarce & — & \zal\ po poprawce D5\\
\bottomrule
\end{tabularx}
\caption{Wyniki testów systemowych}
\end{table}

Weryfikacja ręczna wykryła błąd użyteczności D5: kalendarz wyboru daty zostawał otwarty
po wybraniu dnia i zasłaniał przycisk \enquote{Utwórz}. Po poprawce kalendarz zamyka się
po wyborze, a scenariusz przeszedł w całości, łącznie z utrwaleniem stanu 1/1 po odświeżeniu
strony.

% -------------------------------------------------------------------------
\section{Wyniki testów wydajnościowych}

\begin{table}[H]
\centering\small
\begin{tabular}{@{}l rrr@{}}
\toprule
\textsf{\textbf{Miara}} & \textsf{\textbf{TW-01 smoke}} & \textsf{\textbf{TW-02 load}} & \textsf{\textbf{TW-03 stress}}\\
\midrule
Wirtualni użytkownicy (maks.) & 1 & 10 & 56\\
Iteracje & 38 & 787 & 2\,297\\
Żądania HTTP & 115 & 1\,465 & 4\,043\\
Przepustowość [żądania/s] & 11,25 & 72,66 & 108,59\\
Średni czas [ms] & 22,36 & 21,61 & 250,96\\
Mediana [ms] & 2,18 & 3,84 & 185,08\\
p90 [ms] & 64,00 & 61,34 & 557,87\\
\textbf{p95 [ms]} (próg 300) & \textbf{65,91} & \textbf{64,16} & \textbf{653,32}\\
Maksimum [ms] & 70,58 & 119,08 & 2\,180\\
Błędy HTTP (próg 1\%) & 0,00\% & 0,00\% & 0,00\%\\
Sprawdzenia zaliczone & 114/114 & 1\,464/1\,464 & 4\,042/4\,042\\
\midrule
Wynik & \zal & \zal & ponad próg\\
\bottomrule
\end{tabular}
\caption{Wyniki testów wydajnościowych k6}
\end{table}

\begin{center}
\begin{tikzpicture}
\begin{axis}[
  ybar=3pt, bar width=14pt, width=0.95\textwidth, height=5.8cm,
  ymin=0, ymax=800, ytick={0,200,400,600,800}, ylabel={czas odpowiedzi [ms]},
  symbolic x coords={smoke,load,stress}, xtick=data, enlarge x limits=0.22,
  legend style={at={(0.5,1.04)},anchor=south,legend columns=3,draw=none},
  tick label style={font=\small}, label style={font=\small},
  nodes near coords, every node near coord/.append style={font=\scriptsize},
  nodes near coords style={/pgf/number format/.cd,fixed,precision=0},
  axis line style={draw=muted}, major grid style={draw=muted!25}, ymajorgrids]
\draw[bad,dashed,thick] ({rel axis cs:0,0}|-{axis cs:smoke,300}) -- ({rel axis cs:1,0}|-{axis cs:smoke,300})
  node[pos=0.98,above left,font=\scriptsize,text=bad] {próg p95 = 300 ms};
\addplot[fill=accent!35,draw=accent] coordinates {(smoke,22.36) (load,21.61) (stress,250.96)};
\addplot[fill=accent,draw=accent] coordinates {(smoke,65.91) (load,64.16) (stress,653.32)};
\legend{średnia, p95}
\end{axis}
\end{tikzpicture}
\end{center}

\subsection{Analiza}
\begin{itemize}
  \item Smoke i load zostają pod progiem: p95 wynosi 66 i 64~ms. Zapas względem 300~ms
        jest prawie pięciokrotny przy najwyżej 10 wirtualnych użytkownikach.
  \item Stress narasta do 16 użytkowników logujących się jednocześnie i 40 przeglądających.
        p95 całego przebiegu wynosi 653~ms, mediana 185~ms, maksimum 2,18~s. Próg 300~ms
        jest przekroczony. Błędów HTTP nie ma: ograniczeniem jest czas odpowiedzi, nie dostępność.
  \item Wolne są logowania, bo bcrypt z kosztem 10 blokuje jedyny wątek procesu. Odczyty
        i zmiany statusu czekają w tej samej kolejce, więc przy szczycie zwalnia cały ruch.
  \item Klient k6 i API działały na tej samej maszynie, przez pętlę zwrotną i w trybie
        deweloperskim (\code{tsx}). Wyniki są optymistyczne względem wdrożenia sieciowego.
\end{itemize}

% -------------------------------------------------------------------------
\section{Wyniki audytu bezpieczeństwa}

\subsection{Testy automatyczne i zależności}
Przypadki TB-01–TB-07 zakończyły się wynikiem pozytywnym. \code{pnpm audit} nie zgłosił
żadnej znanej podatności. Wcześniej audyt wykrył podatność krytyczną w pakiecie
\code{tinypool} (zależność Vitest~3) i średnią GHSA-82fw-gwwq-j7x9 w Vitest poniżej 4.1.11.
Obie usunięto przez aktualizację Vitest do 4.1.11 i nadpisanie wersji \code{tinypool}
w \code{pnpm-workspace.yaml} (defekt D4).

\subsection{OWASP ZAP baseline}
Skan ZAP 2.17.0 objął \code{http://127.0.0.1:3000}. Wynik: \textbf{66 reguł zaliczonych,
1 ostrzeżenie, 0 alertów FAIL}.

\begin{table}[H]
\centering\small
\begin{tabularx}{\textwidth}{@{}l L L l@{}}
\toprule
\textsf{\textbf{Reguła}} & \textsf{\textbf{Alert}} & \textsf{\textbf{Adresy}} & \textsf{\textbf{Waga}}\\
\midrule
10049 & Storable and Cacheable Content & \code{/} i \code{/robots.txt} (odpowiedzi 404) & \niski\\
\bottomrule
\end{tabularx}
\caption{Alerty OWASP ZAP}
\end{table}

Pająk ZAP znalazł tylko 3 adresy, bo API nie serwuje strony pod \code{/}, a skan był
nieuwierzytelniony. Wynik potwierdza poprawne nagłówki odpowiedzi publicznych, ale nie
obejmuje tras za logowaniem ani frontendu.

\subsection{Ustalenia przeglądu kodu}
\begin{xltabular}{\textwidth}{@{}>{\sffamily\bfseries\color{accent}}l L l@{}}
\toprule
\textcolor{black}{ID} & \textsf{\textbf{Ustalenie}} & \textsf{\textbf{Waga}}\\
\midrule
\endfirsthead
\toprule
\textcolor{black}{ID} & \textsf{\textbf{Ustalenie}} & \textsf{\textbf{Waga}}\\
\midrule
\endhead
\bottomrule
\endfoot
B1 & \textbf{Klucz limitu logowań pochodzi z nagłówka \code{X-Forwarded-For}}
     (\code{apps/api/src/middleware/rateLimit.ts}). Nagłówek ustawia klient, więc atakujący
     omija limit, zmieniając jego wartość przy każdej próbie. Bez nagłówka wszyscy klienci
     trafiają do wspólnego koszyka \code{local}, więc jeden atakujący może zablokować
     logowanie wszystkim użytkownikom. TI-12 potwierdza, że o kluczu decyduje nagłówek. & \sredni\\
B2 & \textbf{Tokeny sesji są przechowywane jawnie.} Token jest kluczem tabeli \code{sessions}.
     Wyciek pliku bazy pozwala przejąć każdą aktywną sesję na okres do 7 dni. & \sredni\\
B3 & \textbf{Brak \code{Cache-Control: no-store} w odpowiedziach JSON} (alert ZAP 10049).
     Pośrednie bufory mogą przechowywać odpowiedzi API. & \niski\\
B4 & \textbf{Brak rejestrowania zdarzeń bezpieczeństwa} (OWASP A09). Nieudane logowania,
     odpowiedzi 403 i 429 nie trafiają do dziennika, więc atak nie zostawia śladu. & \niski\\
B5 & \textbf{Limit logowań jest trzymany w pamięci procesu.} Resetuje się przy restarcie
     i nie działa wspólnie dla kilku instancji. & \niski\\
B6 & \textbf{Ochrona przed CSRF opiera się wyłącznie na \code{SameSite=Lax}.} API odczytuje
     JSON bez sprawdzania nagłówka \code{Content-Type}. Atak z tej samej witryny
     (np. z subdomeny) nie byłby zablokowany. & \info\\
B7 & \textbf{Ograniczony zasięg DAST}: 3 adresy, bez sesji i bez frontendu. & \info\\
B8 & \textbf{Polityka haseł}: minimum 8 znaków, bcrypt z kosztem 10, bez sprawdzania haseł
     z wycieków. & \info\\
\end{xltabular}

\subsection{Mechanizmy ocenione pozytywnie}
\begin{itemize}
  \item wszystkie zapytania SQL są parametryzowane; TB-04 potwierdza brak wstrzyknięcia;
  \item hasła haszowane bcrypt; ten sam komunikat i porównanie z fikcyjnym skrótem dla
        nieznanego konta ograniczają wyliczanie kont (TI-02);
  \item token sesji ma 256 bitów losowości, ciasteczko \code{HttpOnly}, \code{SameSite=Lax},
        a w produkcji \code{Secure};
  \item cudze zasoby zwracają 404, więc nie da się sprawdzić istnienia listy (TB-06);
  \item przesłane pliki są sprawdzane po typie, rozmiarze i sygnaturze bajtowej (TB-07);
  \item limit ciała żądania 3~MB i ogólny komunikat przy błędzie 500 (TI-13);
  \item lockfile przechodzi politykę łańcucha dostaw pnpm, a audyt zależności jest czysty.
\end{itemize}

% -------------------------------------------------------------------------
\section{Defekty wykryte w trakcie testów}

\begin{table}[H]
\centering\footnotesize
\begin{tabularx}{\textwidth}{@{}>{\sffamily\bfseries\color{accent}}l l L L@{}}
\toprule
\textcolor{black}{ID} & \textsf{\textbf{Wykrył}} & \textsf{\textbf{Opis}} & \textsf{\textbf{Rozwiązanie}}\\
\midrule
D1 & testy integracyjne & błędne ścieżki importu (\code{../} zamiast \code{./}) w \code{context.ts}, \code{http.ts}, \code{uploads.ts}; API nie startowało & poprawione ścieżki\\
D2 & kontrola typów & jednostka pozycji miała typ \code{string} zamiast listy dozwolonych wartości & schemat \code{z.enum(UNITS)}\\
D3 & testy komponentów & DOM przenikał między testami; brak \code{ResizeObserver} w jsdom & sprzątanie po teście i atrapy w \code{setup.ts}\\
D4 & \code{pnpm audit} & podatność krytyczna w \code{tinypool} i średnia w Vitest & Vitest 4.1.11, nadpisanie \code{tinypool}\\
D5 & test ręczny & kalendarz zostawał otwarty i zasłaniał przycisk \enquote{Utwórz} & zamknięcie kalendarza po wyborze dnia\\
D6 & przygotowanie testów k6 & limit 10 logowań na minutę uniemożliwiał pomiar & zmienna \code{LOGIN\_RATE\_LIMIT\_MAX} dla środowiska testowego\\
D7 & przygotowanie skanu ZAP & kontener nie widział API na \code{127.0.0.1} & uruchomienie ZAP w sieci hosta\\
\bottomrule
\end{tabularx}
\caption{Defekty wykryte i usunięte}
\end{table}

% -------------------------------------------------------------------------
\section{Wnioski i zalecenia dla zespołu developerskiego}

\subsection{Wnioski}
\begin{enumerate}
  \item Funkcjonalność jest kompletna i stabilna: przypadki jednostkowe, integracyjne,
        systemowe i bezpieczeństwa przechodzą, a reguła współdzielenia jest potwierdzona
        na poziomie domeny, API i przeglądarki.
  \item Piramida testów działa zgodnie z założeniem: większość defektów (D1–D4) wykryły
        szybkie testy niższych poziomów, a test ręczny wykrył problem użyteczności (D5),
        którego nie opisywał żaden test automatyczny.
  \item Wydajność spełnia wymaganie R8 do 10 wirtualnych użytkowników. Przy 56 użytkownikach,
        w tym 16 logujących się jednocześnie, p95 przekracza 300~ms (653~ms), nadal bez błędów HTTP.
  \item Najistotniejsze ryzyko bezpieczeństwa leży w konfiguracji limitu logowań (B1),
        a nie w obsłudze danych wejściowych. Ta jest solidna.
\end{enumerate}

\subsection{Zalecenia}
\begin{xltabular}{\textwidth}{@{}>{\sffamily\bfseries\color{accent}}l L l@{}}
\toprule
\textcolor{black}{Nr} & \textsf{\textbf{Zalecenie}} & \textsf{\textbf{Priorytet}}\\
\midrule
\endfirsthead
\toprule
\textcolor{black}{Nr} & \textsf{\textbf{Zalecenie}} & \textsf{\textbf{Priorytet}}\\
\midrule
\endhead
\bottomrule
\endfoot
Z1 & Kluczować limit logowań adresem gniazda; \code{X-Forwarded-For} uwzględniać tylko za
     skonfigurowanym, zaufanym proxy. Dodać drugi limit na konto (e-mail). Dopisać test,
     w którym zmiana nagłówka nie omija limitu (B1). & \wysoki\\
Z2 & Zapisywać w bazie skrót SHA-256 tokenu sesji zamiast samego tokenu; przy weryfikacji
     porównywać skróty (B2). & \sredni\\
Z3 & Dodać \code{Cache-Control: no-store} do odpowiedzi JSON API, zostawiając obecne
     ustawienie dla zdjęć (B3). & \sredni\\
Z4 & Rozszerzyć DAST: uwierzytelniony skan ZAP z sesją testową, \code{zap-api-scan} na
     podstawie opisu OpenAPI i skan zbudowanego frontendu (B7). & \sredni\\
Z5 & Uruchamiać testy jednostkowe z progiem pokrycia, integracyjne, bezpieczeństwa i E2E
     w CI przy każdej zmianie; k6 i ZAP cyklicznie. Zapisywać kolejne cykle TDD jako osobne
     commity, aby historia pokazywała przejście od czerwonego do zielonego testu. & \sredni\\
Z6 & Rejestrować nieudane logowania, odpowiedzi 403 i 429 z identyfikatorem konta
     i adresem (B4). & \niski\\
Z7 & W k6 dodać progi per operacja (\code{http\_req\_duration\{name:login\}}),
     test długotrwały (soak) i stress do pierwszych błędów. Mierzyć zbudowane API
     na osobnej maszynie. & \niski\\
Z8 & W testach systemowych dodać Firefox i WebKit oraz scenariusze negatywne: złe hasło,
     wejście obcej osoby na adres cudzej listy. & \niski\\
Z9 & Podnieść pokrycie gałęzi frontu: odmontowanie \code{AuthProvider}, zmiana ilości
     i jednostki w oknie edycji pozycji. Usunąć \code{users.findById} i \code{requireToggle}
     albo użyć ich w kodzie produkcyjnym — dziś wywołują je tylko testy. & \niski\\
Z10 & Przy skalowaniu przenieść limit prób do współdzielonego magazynu (np. Redis);
      rozważyć koszt bcrypt 12 i sprawdzanie haseł z wycieków (B5, B8). & \niski\\
\end{xltabular}

\begin{uwaga}
\small\textbf{Przed wdrożeniem produkcyjnym} należy wykonać co najmniej zalecenia Z1 i Z2.
Pozostałe zalecenia poprawiają odporność i zasięg testów, ale nie blokują wydania.
\end{uwaga}

% -------------------------------------------------------------------------
\appendix
\section{Odtworzenie wyników}

\begin{table}[H]
\centering\small
\begin{tabularx}{\textwidth}{@{}l L@{}}
\toprule
\textsf{\textbf{Polecenie}} & \textsf{\textbf{Działanie}}\\
\midrule
\code{pnpm install} & instalacja zależności\\
\code{pnpm test:unit} & testy jednostkowe i integracyjne z progami pokrycia\\
\code{pnpm test:integration} & tylko testy integracyjne\\
\code{pnpm --filter @shopping/e2e install:browsers} & instalacja Chromium dla Playwright\\
\code{pnpm test:e2e} & testy systemowe\\
\code{pnpm test:perf} & seed i trzy scenariusze k6 na osobnym API\\
\code{pnpm test:security} & testy bezpieczeństwa i \code{pnpm audit}\\
\code{pnpm dev:api}, \code{pnpm test:zap} & API i skan OWASP ZAP baseline\\
\code{tectonic docs/raport/sprawozdanie.tex} & kompilacja tego dokumentu\\
\bottomrule
\end{tabularx}
\caption{Polecenia}
\end{table}

\end{document}
